\documentclass{article}
\usepackage{graphicx} % Required for inserting images

\usepackage{revy}
\usepackage[utf8]{inputenc}
\usepackage[T1]{fontenc}
\usepackage[danish]{babel}
\usepackage{amsmath}
\usepackage{amssymb}   % Required for blackboard bold (\mathbb) and fraktur (\mathfrak)
\usepackage{mathrsfs}


\revyname{MatematikRevyen}
\revyyear{2026}
\version{1}
\eta{$2$ minutter}
\status{Faerdig}

\title{Matematikeren Finn}
\author{Daniel '23, Adam '23, Ronni '25, Peter '23, Louise Carøe '23}

\begin{document}
\maketitle

\begin{roles}
\role{IW}[Interviewer]
\role{F}[Matematikeren Finn]
\end{roles}

\begin{sketch}
\noindent\textit{Skrevet på Gentleman Finn, se inspiration til udførelse her:}\\
www.youtube.com/watch?v=VPS8GouqQgs\\

*Eventuelt fisk slutter*
*Lys op*

\says{IW} Og nu videre til noget HELT andet. Det er jo en matematikrevy, hvor matematikken er i centrum. Derfor vil jeg gerne byde velkommen til en person der har været på en lang rejse for at finde sit lod i livet. Jeg byder velkommen til dig, Matematikeren Finn!\\

\scene{F kommer ind på scenen og sætter sig ved bordet ved siden af IW}\\

\says{F} Mange tak skal du have IW *Mens han giver IW hånden*

\says{IW} Finn, du har jo både været satanist, anti-vaxxer, flat earther og nazist, men det har du sagt farvel til.

\says{F} Au revior, IW. Jeg har jo flere gange kommet ud på et sidespor, hvor du har hjulpet mig, og er jeg meget taknemmelig for din hjælp.

\says{IW} Sådan, Finn. Hvor er det godt at høre, at du har kommet væk fra dine afveje.

\says{F} Jaja, og jeg har slet ikke tid til det der andet pjat, da jeg er blevet... forsikringsmatematiker.

\says{IW} ...Eeej, for helvede, Finn. Det er da det værste, du kunne være. Mener du det?

\says{F} Jojo, vi er nogle friske unge fyre, som godt kan lide at tage ind på HCØ og beregne på, hvornår vores forældre dør.

\says{IW} Ej, det mener du altså ikke? Det er da fuldstændig amoralsk.

\says{F} Neearj, det er jo et projekt. Endda statsstøttet og fint beskrevet på flere af KU's hjemmesider.

\says{IW} Men det gør det da ikke moralsk forsvarligt. Tag eksempelvis Niels Bohr bygningen, den har jo også været statsstøttet.

\says{F} OOkkay, den havde jeg ikke lige tænkt på, point taken, point taken. Men jeg er stadig matematiker, og du skal jo tænke på, at os forsikringsmatematikere har givet jer en masse gode ting.

\says{IW} Som hvad?

\says{F} Vi har jo forsikret jer om, at der er uendelige mange primtal.

\says{IW}. Men men, det er jo abstrakt matematisk deduktion.

\says{F} ...og?

\says{IW} Det har jo ikke noget med forsikringsmatematik at gøre.

\says{F} jojo

\says{IW} Nej altså forsikringer er jo noget du skal have hvis du skal rejse eller have et hus eller sådan noget

\says{F} Når når OOOOOKKKKAY, ja, ja, der twister du den. Sådan havde jeg ikke lige set den-

\textbf{IW}: Ja, man skal se den fra flere sider-

\says{F} Jaja, du taler godt for dig. Det tør siges.

\says{IW} Men altså det betyder, at du er jo ikke rigtig matematiker.

\says{F} Okay, så måske bare en halv matematiker, fordi jeg kommer stadig ude på Science.

\says{IW} Jaa, men du beregner på, hvor dine forældre dør. Altså du er ulækker også. Den skal du lige have med.

\says{F} Jaja, okay, point taken, du twister den fint, så måske bare en ulækker kvart matematiker. Meeeen du skal også lige tænke på her i det store regnestykke, at vi regner også, hvornår DINE forældre dør

\says{IW} Det gør det jo ikke meget bedre, det skal du altså lade være med, jeg gider da ikke vide, hvornår de dør. 

\says{F} Men det er jo DÈR pengene ligger.

\says{IW} Nej, men matematik handler jo om skønheden, logikken, elegancen ved det hele, ikke om penge.

\says{F} Ooookay, dér er du hurtig. Det er du. Jaja, men du skal tænke på, at når jeg bruger mit lille kridt på tavlen, finder ud af hvornår mine forældre dør og får stukket penge i hånden, så var det jo af kærlighed og ikke grådighed.

\says{IW}\noindent\act{krydser armen} Nej, det er jo LIGE omvendt.

\says{F} HVA?! ER  DER NOGLE, SOM STIKKER DERES LILLE KRIDT IND I EN TAVLE OG FÅR NOGLE FORÆLDRE UD AF DET???!

\says{IW} Finn prøv at høre, du skal stoppe med at det der gravkammer-statistik, du er heller ikke matematiker. I så fald skal vi ned til en ottenedel.

\says{F} Ja, ja, ja, du taler godt for dig... Nej, men du skal lige huske, at nogle af de penge jeg tjener på at regne på deres død, bruger jeg på deres begravelse, så en møg ulækker sjettedels matematiker, og så skal jeg ikke længere ned

\says{IW} Men må jeg lige sige noget?  Du køber jo hverken blomster eller gravøl, så du er jo skide ligeglad med dine forældre. Så mit bud er hensynsløs, ulækker OG SÅ en syvendedels matematiker

\noindent\act{IW stikker hånden ud}

\says{F} ...Så bliver jeg jo nødt til at sige top

\noindent\act{Giver IW hånden}

\says{IW + F} Det er en aftale.

\noindent\act{IW kigger ned i sine papirer og tager det næste frem}

\says{IW} Så en syvendedel blev vi enige om...
Og så er du jo også fysiker ved siden af, men det talte jeg ikke med.

\says{F} Nej nej, for det har ikke noget sagen at gøre.

\says{IW} Overhovedet ikke, en gang matematiker, altid matematiker

\end{sketch}
\end{document}
